\section{Introduction}
% GOAL: Position as evaluation FRAMEWORK, not benchmark.
% Structure: Problem → Insight → Framework Design → Contributions

% === TENSION STRUCTURE ===
% 1. Agentic models are moving to production systems → gap in evaluation
% 2. Existing benchmarks are static, isolated, short-horizon → problem
% 3. Real production engineering is serial, stateful, multi-stage → insight
% 4. We design an evaluation framework (not a benchmark) → approach
% 5. Contributions: Framework + Workloads + Metrics + Results → closure

Large language model (LLM) agents are rapidly evolving from code assistants into autonomous software engineering systems capable of planning, implementing, debugging, and deploying complex infrastructure stacks~\citep{jimenez2024swebench,liu2024agentbench,frontiereng2026}.
Yet the evaluation methodologies for these systems remain rooted in a \emph{benchmark-driven paradigm}: isolated, static, short-horizon tasks where each evaluation unit is independent and the agent is reset after every attempt.

This paradigm has a fundamental limitation: \emph{real production engineering is not a collection of independent tasks}.
Software development proceeds along dependency chains---a compiler's parser cannot function without a working lexer; a deployment pipeline cannot succeed without correct build artifacts.
When an agent fails at an early stage, the entire downstream workflow becomes unreachable, not because the agent lacks downstream capability, but because the prerequisite artifact is missing.
Static benchmarks cannot distinguish between these two failure modes, conflating ``cannot reach'' with ``cannot solve.''

\textbf{Insight.}
We argue that the field needs to move \emph{beyond benchmarks} toward \emph{evaluation frameworks}---systems that assess agentic models in realistic, multi-stage production environments with controlled intermediate state injection, multi-dimensional metrics, and systematic failure diagnosis.
A benchmark provides a fixed dataset and scoring function; an evaluation framework provides a \emph{methodology} for constructing workloads, injecting controlled interventions, and measuring both outcomes and processes.

\textbf{Approach.}
We present \evobench{}, an evaluation framework for assessing agentic models in production-grounded software engineering environments.
\evobench{} is built on four pillars:
\begin{enumerate}[nosep]
\item \textbf{Framework}: A unified evaluation infrastructure (\textsc{AIhub}) that supports multi-model access (21 models) and multi-SDK integration (OpenHands, Claude Code, Codex CLI, Kimi CLI), with containerized isolation for reproducible execution.
\item \textbf{Workloads}: Two complementary compiler construction workloads---\yatcc{} (scaffolded fill-in-the-blank tasks) and \yatcc{}-Hard (full implementation from scratch with no code logic provided)---both organized as serial dependency chains where each task's output becomes the next task's input.
\item \textbf{Metrics}: A multi-dimensional evaluation system spanning outcome metrics (Mean Reward, Pass Score, Pipeline Score), process metrics (tokens, turns, retries, elapsed time), and a five-category failure taxonomy (predecessor failure, task-local reasoning failure, execution failure, process collapse, cost overrun).
\item \textbf{Resurrection}: A controlled intervention mechanism that injects verified intermediate artifacts after task failure, decomposing cascading failure into predecessor-unreachable and downstream-incapable, enabling node-level diagnostic evaluation even when the full pipeline is broken.
\end{enumerate}

Our contributions are as follows:
\begin{enumerate}[nosep]
\item We propose \evobench{} as an \emph{evaluation framework} (not a benchmark) for production-grounded assessment of agentic models, with unified multi-model/SDK access, containerized isolation, and reproducible execution.
\item We design two complementary \emph{serial workloads}---\yatcc{} and \yatcc{}-Hard---that form causal dependency chains with verifiable intermediate states, spanning a difficulty spectrum from scaffolded implementation to full from-scratch construction.
\item We introduce \resurrection{} as an \emph{evaluation methodology} that decomposes cascading failure into diagnostic components, recovering downstream evaluation signal that would otherwise be lost (median 73.65-point pipeline score recovery across 19 configurations).
\item We establish a \emph{multi-dimensional metrics system} where process cost (tokens, turns, retries) is a first-class analytical object alongside outcome scores, revealing that process cost is decoupled from outcome ($R^2 = 0.007$)---a finding with direct implications for production deployment decisions.
\item We provide empirical evidence across 22 model configurations, four agent backends, and two workload difficulty levels, demonstrating that serial evaluation preserves overall rankings ($\tau = 0.763$) but reveals specific failure patterns invisible to static benchmarks, and that the difficulty gap between \yatcc{} and \yatcc{}-Hard quantifies the value of scaffolding in agentic coding tasks.
\end{enumerate}